% !TeX TS-program = pdflatex


\documentclass[a4paper]{article}

% \usepackage[default]{fontsetup}

\usepackage{fancyhdr}
\usepackage{extramarks}
\usepackage{amsmath}
\usepackage{amsthm}
\usepackage{amsfonts}
\usepackage{tikz}
\usepackage[plain]{algorithm}
\usepackage{algpseudocode}
\usepackage{enumerate}
\usepackage{tikz}

\usetikzlibrary{automata,positioning}

%
% Basic Document Settings
%  

\topmargin=-0.2in
\evensidemargin=0in
\oddsidemargin=0in
\textwidth=6.5in
\textheight=9.5in
\headsep=0.25in

\linespread{1.1}

\pagestyle{fancy}
\lhead{\hmwkAuthorName}
\chead{\hmwkClass : \hmwkTitle}
\rhead{\firstxmark}
\lfoot{\lastxmark}
\cfoot{\thepage}

\renewcommand\headrulewidth{0.4pt}
\renewcommand\footrulewidth{0.4pt}

\setlength\parindent{0pt}

%
% Create Problem Sections
%

\newcommand{\enterProblemHeader}[1]{
    \nobreak\extramarks{}{Problem \arabic{#1} continued on next page\ldots}\nobreak{}
    \nobreak\extramarks{Problem \arabic{#1} (continued)}{Problem \arabic{#1} continued on next page\ldots}\nobreak{}
}

\newcommand{\exitProblemHeader}[1]{
    \nobreak\extramarks{Problem \arabic{#1} (continued)}{Problem \arabic{#1} continued on next page\ldots}\nobreak{}
    \stepcounter{#1}
    \nobreak\extramarks{Problem \arabic{#1}}{}\nobreak{}
}

\newcommand*\circled[1]{\tikz[baseline=(char.base)]{
		\node[shape=circle,draw,inner sep=2pt] (char) {#1};}}


\setcounter{secnumdepth}{0}
\newcounter{partCounter}
\newcounter{homeworkProblemCounter}
\setcounter{homeworkProblemCounter}{1}
\nobreak\extramarks{Problem \arabic{homeworkProblemCounter}}{}\nobreak{}

%
% Homework Problem Environment
%
% This environment takes an optional argument. When given, it will adjust the
% problem counter. This is useful for when the problems given for your
% assignment aren't sequential. See the last 3 problems of this template for an
% example.
%

\newenvironment{homeworkProblem}[1][-1]{
    \ifnum#1>0
        \setcounter{homeworkProblemCounter}{#1}
    \fi
    \section{Problem \arabic{homeworkProblemCounter}}
    \setcounter{partCounter}{1}
    \enterProblemHeader{homeworkProblemCounter}
}{
    \exitProblemHeader{homeworkProblemCounter}
}

%
% Homework Details
%   - Title
%   - Class
%   - Due date
%   - Name
%   - Student ID

\newcommand{\hmwkTitle}{Homework\ \#02}
\newcommand{\hmwkClass}{Probability \& Statistics for EECS}
\newcommand{\hmwkDueDate}{2024-10-6}
\newcommand{\hmwkAuthorName}{}
\newcommand{\hmwkAuthorID}{}


%
% Title Page
%

\title{
    \vspace{2in}
    \textmd{\textbf{\hmwkClass:\\  \hmwkTitle}}\\
    \normalsize\vspace{0.1in}\small{Due\ on\ \hmwkDueDate\ at 18:00}\\
	\vspace{4in}
}

\author{
	Name: \textbf{\hmwkAuthorName} \\
	Student ID: \hmwkAuthorID}
\date{}

\renewcommand{\part}[1]{\textbf{\large Part \Alph{partCounter}}\stepcounter{partCounter}\\}

%
% Various Helper Commands
%

% Useful for algorithms
\newcommand{\alg}[1]{\textsc{\bfseries \footnotesize #1}}
% For derivatives
\newcommand{\deriv}[1]{\frac{\mathrm{d}}{\mathrm{d}x} (#1)}
% For partial derivatives
\newcommand{\pderiv}[2]{\frac{\partial}{\partial #1} (#2)}
% Integral dx
\newcommand{\dx}{\mathrm{d}x}
% Alias for the Solution section header
\newcommand{\solution}{\textbf{\large Solution}}
% Probability commands: Expectation, Variance, Covariance, Bias
\newcommand{\E}{\mathrm{E}}
\newcommand{\Var}{\mathrm{Var}}
\newcommand{\Cov}{\mathrm{Cov}}
\newcommand{\Bias}{\mathrm{Bias}}

\begin{document}


% \maketitle
% \thispagestyle{empty}
% \pagebreak

\date{
Due on Oct. 7, 2026, 23:59 UTC+8}
\title{SI 140A-01  Probability \& Statistics for EECS, Fall 2026 \\
Homework 2}
\maketitle
Read all the instructions below carefully before you start working on the assignment, and before you make a submission.
\begin{itemize}
    \item You are required to write down all the major steps towards making your conclusions; otherwise you may obtain limited points of the problem.
    \item Write your homework in English; otherwise you will get no points of this homework.
    \item Any form of plagiarism will lead to $0$ point of this homework. 
\end{itemize}
\newpage

\begin{homeworkProblem}[1]
Show that for any events $A$ and $B$,
    \[
    P(A) + P(B) - 1 \le P(A \cap B) \le P(A \cup B) \le P(A) + P(B).
    \]
    For each of these three inequalities, give a simple criterion for when the inequality is actually an equality (e.g., give a simple condition such that $P(A \cap B) = P(A \cup B)$ if and only if the condition holds).
\end{homeworkProblem}

\newpage
\begin{homeworkProblem}[2]
Let $A$ and $B$ be events. The \emph{symmetric difference} $A \triangle B$ is defined to be the set of all elements that are in $A$ or $B$ but not both. In logic and engineering, this event is also called the $\text{XOR}$ (\emph{exclusive or}) of $A$ and $B$. Show that
    \[
    P(A \triangle B) = P(A) + P(B) - 2P(A \cap B),
    \]
    directly using the axioms of probability.
\end{homeworkProblem}



%\newpage
%\begin{homeworkProblem}[3]
%Given $n \geq 2$ different numbers $\left(a_1, a_2, \ldots, a_n\right)$, a bootstrap sample is a sequence $\left(x_1, x_2, \ldots, x_n\right)$ formed from the $a_j$ 's by sampling with replacement with equal probabilities.  For example, if $n=2$ and $\left(a_1, a_2\right)=(3,1)$, then the possible bootstrap samples are $(3,3),(3,1),(1,3)$, and $(1,1)$.\\
%(a) How many possible bootstrap samples are there for $\left(a_1, \ldots, a_n\right)$ ?\\
%(b) How many possible bootstrap samples are there for $\left(a_1, \ldots, a_n\right)$, if order does not matter.\\ %(in the sense that it only matters how many times each $a_j$ was chosen, not the order in which they were chosen)?\\
%(c) One random bootstrap sample is chosen (by sampling from $a_1, \ldots, a_n$ with replacement, as described above). Show that not all unordered bootstrap samples (in the sense of (b)) are equally likely. Find an unordered bootstrap sample $\mathbf{b}_1$ that is as likely as possible, and an unordered bootstrap sample $\mathbf{b}_2$ that is as unlikely as possible. Let $p_1$ be the probability of getting $\mathbf{b}_1$ and $p_2$ be the probability of getting $\mathbf{b}_2$ (so $p_i$ is the probability of getting the specific unordered bootstrap sample $\mathbf{b}_i$ ). What is $p_1 / p_2$ ? What is the ratio of the probability of getting an unordered bootstrap sample whose probability is $p_1$ to the probability of getting an unordered sample whose probability is $p_2$ ?

%\end{homeworkProblem}

% \newpage
% \begin{homeworkProblem}[4]
% \textbf{(Geometric Probability)} You get a stick and break it randomly into three pieces. What is the probability that you can make a triangle using such three pieces?
% \end{homeworkProblem}


\newpage
\begin{homeworkProblem}[3]
Events $A$ and $B$ are \emph{independent} if $P(A \cap B) = P(A)P(B)$.
    \begin{enumerate}
        \item Give an example of independent events $A$ and $B$ in a finite sample space $S$ (with neither equal to $\emptyset$ or $S$).
        \item Consider the experiment of picking a random point in the rectangle
        \[
        R = \{(x, y) : 0 < x < 1, 0 < y < 1\},
        \]
        where the probability of the point being in any particular region contained within $R$ is the area of that region. Let $A_1$ and $B_1$ be rectangles contained within $R$, with areas not equal to 0 or 1. Let $A$ be the event that the random point is in $A_1$, and $B$ be the event that the random point is in $B_1$. Give a geometric description of when it is true that $A$ and $B$ are independent. Also, give an example where they are independent and another example where they are not independent.
        \item Show that if $A$ and $B$ are independent, then
        \[
        P(A \cup B) = P(A) + P(B) - P(A)P(B) = 1 - P(A^c)P(B^c).
        \]
    \end{enumerate}

\end{homeworkProblem}

% \newpage
% \begin{homeworkProblem}[4]
% \textbf{(Geometric Probability)} You get a stick and break it randomly into three pieces. What is the probability that you can make a triangle using such three pieces?
% \end{homeworkProblem}

\newpage
\begin{homeworkProblem}[4]
Consider four nonstandard dice (the \emph{Efron dice}), whose sides are labeled as follows (the 6 sides on each die are equally likely).
    \begin{itemize}
        \item[] A: $4, 4, 4, 4, 0, 0$
        \item[] B: $3, 3, 3, 3, 3, 3$
        \item[] C: $6, 6, 2, 2, 2, 2$
        \item[] D: $5, 5, 5, 1, 1, 1$
    \end{itemize}
    These four dice are each rolled once. Let $A$ be the result for die $A$, $B$ be the result for die $B$, etc.
    \begin{enumerate}
        \item Find $P(A > B)$, $P(B > C)$, $P(C > D)$, and $P(D > A)$.
        \item Is the event $A > B$ independent of the event $B > C$? Is the event $B > C$ independent of the event $C > D$? Explain.
    \end{enumerate}
\end{homeworkProblem}

 




\end{document}



% % !TeX TS-program = pdflatex


% \documentclass[a4paper]{article}

% % \usepackage[default]{fontsetup}

% \usepackage{fancyhdr}
% \usepackage{extramarks}
% \usepackage{amsmath}
% \usepackage{amsthm}
% \usepackage{amsfonts}
% \usepackage{tikz}
% \usepackage[plain]{algorithm}
% \usepackage{algpseudocode}
% \usepackage{enumerate}
% \usepackage{tikz}

% \usetikzlibrary{automata,positioning}

% %
% % Basic Document Settings
% %  

% \topmargin=-0.2in
% \evensidemargin=0in
% \oddsidemargin=0in
% \textwidth=6.5in
% \textheight=9.5in
% \headsep=0.25in

% \linespread{1.1}

% \pagestyle{fancy}
% \lhead{\hmwkAuthorName}
% \chead{\hmwkClass : \hmwkTitle}
% \rhead{\firstxmark}
% \lfoot{\lastxmark}
% \cfoot{\thepage}

% \renewcommand\headrulewidth{0.4pt}
% \renewcommand\footrulewidth{0.4pt}

% \setlength\parindent{0pt}

% %
% % Create Problem Sections
% %

% \newcommand{\enterProblemHeader}[1]{
%     \nobreak\extramarks{}{Problem \arabic{#1} continued on next page\ldots}\nobreak{}
%     \nobreak\extramarks{Problem \arabic{#1} (continued)}{Problem \arabic{#1} continued on next page\ldots}\nobreak{}
% }

% \newcommand{\exitProblemHeader}[1]{
%     \nobreak\extramarks{Problem \arabic{#1} (continued)}{Problem \arabic{#1} continued on next page\ldots}\nobreak{}
%     \stepcounter{#1}
%     \nobreak\extramarks{Problem \arabic{#1}}{}\nobreak{}
% }

% \newcommand*\circled[1]{\tikz[baseline=(char.base)]{
% 		\node[shape=circle,draw,inner sep=2pt] (char) {#1};}}


% \setcounter{secnumdepth}{0}
% \newcounter{partCounter}
% \newcounter{homeworkProblemCounter}
% \setcounter{homeworkProblemCounter}{1}
% \nobreak\extramarks{Problem \arabic{homeworkProblemCounter}}{}\nobreak{}

% %
% % Homework Problem Environment
% %
% % This environment takes an optional argument. When given, it will adjust the
% % problem counter. This is useful for when the problems given for your
% % assignment aren't sequential. See the last 3 problems of this template for an
% % example.
% %

% \newenvironment{homeworkProblem}[1][-1]{
%     \ifnum#1>0
%         \setcounter{homeworkProblemCounter}{#1}
%     \fi
%     \section{Problem \arabic{homeworkProblemCounter}}
%     \setcounter{partCounter}{1}
%     \enterProblemHeader{homeworkProblemCounter}
% }{
%     \exitProblemHeader{homeworkProblemCounter}
% }

% %
% % Homework Details
% %   - Title
% %   - Class
% %   - Due date
% %   - Name
% %   - Student ID

% \newcommand{\hmwkTitle}{Homework\ \#13}
% \newcommand{\hmwkClass}{Probability \& Statistics for EECS}
% \newcommand{\hmwkDueDate}{2025-01-07}
% \newcommand{\hmwkAuthorName}{}
% \newcommand{\hmwkAuthorID}{}


% %
% % Title Page
% %

% \title{
%     \vspace{2in}
%     \textmd{\textbf{\hmwkClass:\\  \hmwkTitle}}\\
%     \normalsize\vspace{0.1in}\small{Due\ on\ \hmwkDueDate\ at 11:59}\\
% 	\vspace{4in}
% }

% \author{
% 	Name: \textbf{\hmwkAuthorName} \\
% 	Student ID: \hmwkAuthorID}
% \date{}

% \renewcommand{\part}[1]{\textbf{\large Part \Alph{partCounter}}\stepcounter{partCounter}\\}

% %
% % Various Helper Commands
% %

% % Useful for algorithms
% \newcommand{\alg}[1]{\textsc{\bfseries \footnotesize #1}}
% % For derivatives
% \newcommand{\deriv}[1]{\frac{\mathrm{d}}{\mathrm{d}x} (#1)}
% % For partial derivatives
% \newcommand{\pderiv}[2]{\frac{\partial}{\partial #1} (#2)}
% % Integral dx
% \newcommand{\dx}{\mathrm{d}x}
% % Alias for the Solution section header
% \newcommand{\solution}{\textbf{\large Solution}}
% % Probability commands: Expectation, Variance, Covariance, Bias
% \newcommand{\E}{\mathrm{E}}
% \newcommand{\Var}{\mathrm{Var}}
% \newcommand{\Cov}{\mathrm{Cov}}
% \newcommand{\Bias}{\mathrm{Bias}}

% \begin{document}


% % \maketitle
% % \thispagestyle{empty}
% % \pagebreak

% \date{
% Due on Jan. 7, 2025, 11:59 UTC+8}
% \title{SI 140A-02  Probability \& Statistics for EECS, Fall 2024 \\
% Homework 13}
% \maketitle
% Read all the instructions below carefully before you start working on the assignment, and before you make a submission.
% \begin{itemize}
%     \item You are required to write down all the major steps towards making your conclusions; otherwise you may obtain limited points of the problem.
%     \item Write your homework in English; otherwise you will get no points of this homework.
%     \item Any form of plagiarism will lead to $0$ point of this homework. 
% \end{itemize}
% \newpage

% \begin{homeworkProblem}[1]
% Given a Markov chain with state-transition diagram shown as follows:
% \begin{center}
%     \includegraphics[width=0.3\textwidth]{hw13_1.png}
% \end{center}
% (a)
% Is this chain irreducible?\\[2px]
% (b)
% Is this chain aperiodic?\\[2px]
% (c)
% Find the stationary distribution of this chain.\\[2px]
% (d)
% Is this chain reversible?
% \end{homeworkProblem}

% \newpage
% \begin{homeworkProblem}[2]
% Given a Markov chain with state-transition diagram shown as follows:
% \begin{center}
%     \includegraphics[width=0.3\textwidth]{hw13_1.png}
% \end{center}
% (a) Find \(P(X_{9}=3 \mid X_{8}=1)\) and \(P(X_{8}=2 \mid X_{7}=3)\).\\[2px]
% (b) If \(P(X_{0}=3) = \frac{1}{2}\), find \(P(X_{0}=3, X_{1}=1, X_{2}=2, X_{4}=3)\).\\[2px]
% (c) Find \(E(X_{8} \mid X_{6}=2)\).\\[2px]
% (d) Find \(\operatorname{Var}(X_{7} \mid X_{5}=3)\).\\[2px]
% \end{homeworkProblem}



% \newpage
% \begin{homeworkProblem}[3]
% There are two urns with a total of \(2N\) distinguishable balls. Initially, the first urn has \(N\) white balls and the second urn has \(N\) black balls. At each stage, we pick a ball at random from each urn and interchange them. Let \(X_{n}\) be the number of black balls in the first urn at time \(n\). This is a Markov chain on the state space \(\{0,1,\ldots,N\}\).\\[4px]
% (a) Find the transition probabilities of the chain.\\[2px]
% (b) Find the stationary distribution of the chain.\\[2px]

% \end{homeworkProblem}


% \end{document}
